\documentclass[11pt]{article}
% ===========================================================================
%  Shared preamble for the AP Calculus AB exam-preparation worksheets.
%  Input right after \documentclass. Files beginning with "_" are not compiled
%  on their own. Navy and gold, to match the course's lesson pages.
% ===========================================================================
\usepackage[letterpaper,top=2.3cm,bottom=2.2cm,left=1.7cm,right=1.7cm,headheight=28pt,footskip=30pt]{geometry}
\usepackage{amsmath,amssymb}
\usepackage[T1]{fontenc}
\usepackage[utf8]{inputenc}
\usepackage{lmodern}
\usepackage{xcolor}
\usepackage[most]{tcolorbox}
\usepackage{enumitem}
\usepackage{fancyhdr}
\usepackage{array}
\usepackage{tikz}
\usepackage{pgfplots}
\pgfplotsset{compat=1.18}

\definecolor{navy}{HTML}{161E54}
\definecolor{ink}{HTML}{27357E}
\definecolor{gold}{HTML}{B98A00}
\definecolor{goldbg}{HTML}{FFFDF3}
\definecolor{inkbg}{HTML}{F6F7FD}
\definecolor{rule}{HTML}{C9CFF0}
\definecolor{okgreen}{HTML}{1B7A44}
\definecolor{okbg}{HTML}{F1F8F3}

\pagestyle{fancy}
\fancyhf{}
\fancyhead[L]{\footnotesize Name: \underline{\hspace{4cm}}}
\fancyhead[C]{\textbf{\color{navy}Integration Academy}\\[-2pt]{\scriptsize Calculus AB \textperiodcentered{} AP\textsuperscript{\textregistered} Exam Preparation}}
\fancyfoot[L]{\tiny\color{black!55}\parbox[t]{0.78\textwidth}{AP\textsuperscript{\textregistered} is a registered trademark of the College Board, which was not involved in the production of, and does not endorse, this material. \textcopyright{} Integration Academy.}}
\fancyfoot[R]{\scriptsize\thepage}
\renewcommand{\headrulewidth}{1.2pt}
\renewcommand{\headrule}{\hbox to\headwidth{\color{gold}\leaders\hrule height \headrulewidth\hfill}}
\renewcommand{\footrulewidth}{0pt}
\setlength{\parindent}{0pt}
\setlength{\parskip}{4pt}

% Title block: \wstitle{code}{title}{unit line}
\newcommand{\wstitle}[3]{\begin{center}{\small\bfseries\color{gold}WORKSHEET #1}\\[3pt]{\LARGE\bfseries\color{navy}#2}\\[3pt]{\small\color{black!60}#3}\end{center}\vspace{2pt}}
% Section heading
\newcommand{\sech}[1]{\par\vspace{10pt}{\large\bfseries\color{navy}#1}\par\nobreak\vspace{2pt}{\color{rule}\hrule height0.8pt}\vspace{7pt}}
% A level tag printed beside a group of questions
\newcommand{\level}[1]{\par\vspace{6pt}{\footnotesize\bfseries\color{gold}\MakeUppercase{#1}}\par\nobreak\vspace{2pt}}

\newtcolorbox{ideabox}{colframe=ink,colback=inkbg,coltitle=white,colbacktitle=ink,fonttitle=\bfseries\small,title={Key ideas},arc=2pt,boxrule=0.8pt,left=8pt,right=8pt}
\newtcolorbox{example}[1]{colframe=rule,colback=white,coltitle=navy,colbacktitle=inkbg,fonttitle=\bfseries\small,title={#1},arc=2pt,boxrule=0.8pt,borderline west={3pt}{0pt}{ink},left=10pt,right=8pt}
\newtcolorbox{frbox}[1]{colframe=gold,colback=goldbg,coltitle=white,colbacktitle=gold,fonttitle=\bfseries\small,title={#1},arc=2pt,boxrule=0.8pt,breakable,left=8pt,right=8pt}
\newtcolorbox{solbox}{colframe=okgreen,colback=okbg,coltitle=white,colbacktitle=okgreen,fonttitle=\bfseries\small,title={Worked solutions},arc=2pt,boxrule=0.8pt,breakable,left=8pt,right=8pt}
\newtcolorbox{scorebox}[1]{colframe=navy,colback=white,coltitle=white,colbacktitle=navy,fonttitle=\bfseries\small,title={#1},arc=2pt,boxrule=0.8pt,breakable,left=8pt,right=8pt}

\newcommand{\soln}{\par\smallskip\textbf{\color{ink}Solution.}\ }
\newcommand{\result}[1]{\par\smallskip\textbf{#1}}
% blank working space of a given height
\newcommand{\work}[1]{\par\vspace{#1}}
% a point value in the margin of a scoring guide
\newcommand{\pts}[1]{\hfill{\footnotesize\bfseries\color{gold}[#1]}}

% Coordinate grid for graph questions: \begin{apgraph}{xmin}{xmax}{ymin}{ymax} ... \end{apgraph}
\newenvironment{apgraph}[4]{\begin{center}\begin{tikzpicture}\begin{axis}[
  axis lines=middle, xmin=#1, xmax=#2, ymin=#3, ymax=#4,
  grid=both, grid style={black!12}, xtick distance=1, ytick distance=1,
  ticklabel style={font=\scriptsize}, xlabel={$x$}, ylabel={$y$},
  width=10cm, height=7.2cm, samples=80, clip=false,
  every axis plot/.append style={line width=1.1pt, color=ink}]}{\end{axis}\end{tikzpicture}\end{center}}
\newcommand{\closeddot}[2]{\addplot[only marks, mark=*, mark size=2.2pt, color=ink] coordinates {(#1,#2)};}
\newcommand{\opendot}[2]{\addplot[only marks, mark=*, mark size=2.2pt, color=ink, mark options={fill=white}] coordinates {(#1,#2)};}

% Multiple-choice options: four across, or stacked when the options are long.
\newcommand{\opts}[4]{\par\vspace{3pt}\begin{tabular}{@{}p{0.235\linewidth}p{0.235\linewidth}p{0.235\linewidth}p{0.235\linewidth}@{}}(A)\ #1 & (B)\ #2 & (C)\ #3 & (D)\ #4\end{tabular}\par\vspace{9pt}}
\newcommand{\optsv}[4]{\par\vspace{3pt}\begin{tabular}{@{}l@{\ \ }p{0.9\linewidth}@{}}(A) & #1\\(B) & #2\\(C) & #3\\(D) & #4\end{tabular}\par\vspace{9pt}}

\fancyhead[R]{\footnotesize\color{black!60}Worksheet 2.6}
\begin{document}

\wstitle{2.6}{Derivatives of Sine, Cosine, $e^{x}$ and $\ln x$}{Unit 2: The Derivative}

\begin{ideabox}
\begin{itemize}[leftmargin=*,itemsep=2pt,topsep=2pt]
  \item $\dfrac{d}{dx}[\sin x]=\cos x$ \quad and \quad $\dfrac{d}{dx}[\cos x]=-\sin x$ \quad ($x$ in radians).
  \item $\dfrac{d}{dx}\left[e^{x}\right]=e^{x}$ \quad and \quad $\dfrac{d}{dx}[\ln x]=\dfrac1x$ for $x>0$.
  \item These combine with the sum and constant-multiple rules exactly as powers do.
  \item Exact trigonometric values are expected: know sine and cosine at $0,\ \tfrac\pi6,\ \tfrac\pi4,\ \tfrac\pi3,\ \tfrac\pi2$.
\end{itemize}
\end{ideabox}

\sech{Worked examples}
\begin{example}{Example 1 \textperiodcentered{} A sum}
Find $\dfrac{d}{dx}\left[3\sin x-4\cos x+2e^{x}\right]$.\soln Differentiate each term; the derivative of $-4\cos x$ is $+4\sin x$.\result{$3\cos x+4\sin x+2e^{x}$.}
\end{example}

\begin{example}{Example 2 \textperiodcentered{} A tangent line}
Write an equation of the line tangent to $y=\ln x$ at $x=2$.\soln The point is $(2,\ln2)$ and the slope is $\dfrac12$.\result{$y-\ln2=\dfrac12(x-2)$.}
\end{example}

\begin{example}{Example 3 \textperiodcentered{} Horizontal tangents}
Find the values of $x$ in $[0,2\pi]$ at which $y=\sqrt3\,x+2\cos x$ has a horizontal tangent.\soln $y'=\sqrt3-2\sin x=0$ gives $\sin x=\dfrac{\sqrt3}{2}$.\result{$x=\dfrac\pi3$ and $x=\dfrac{2\pi}{3}$.}
\end{example}

\newpage
\sech{Practice}

\level{Core}
\begin{enumerate}[leftmargin=*,itemsep=2pt]
  \item Find $\dfrac{d}{dx}\left[5e^{x}\right]$.\work{2.3cm}
  \item Find $\dfrac{d}{dx}\left[2\sin x+x^3\right]$.\work{2.3cm}
  \item Find $\dfrac{d}{dx}\left[4\ln x-x\right]$.\work{2.3cm}
  \item For $f(x)=\cos x$, find $f'\!\left(\dfrac\pi6\right)$.\work{2.3cm}
\end{enumerate}

\level{Exam level}
\begin{enumerate}[leftmargin=*,itemsep=2pt,start=5]
  \item Write an equation of the line tangent to $y=2\sin x$ at $x=\dfrac\pi3$.\work{2.8cm}
  \item At what value of $x$ does the graph of $y=2x-e^{x}$ have a horizontal tangent?\work{2.8cm}
  \item Evaluate $\displaystyle\lim_{h\to 0}\dfrac{e^{3+h}-e^{3}}{h}$.\work{2.8cm}
  \item Let $f(x)=6\ln x+x^2$ for $x>0$. Find all values of $x$ at which $f'(x)=7$.\work{2.8cm}
\end{enumerate}

\begin{frbox}{Question 9 \textperiodcentered{} Free-response practice}
Let $f(x)=4\sin x+2x$ for $0\le x\le2\pi$.
\begin{enumerate}[label=(\alph*),leftmargin=*,itemsep=2pt]
  \item Find $f'(x)$.\work{2.6cm}
  \item Find all values of $x$ in the interval at which the graph of $f$ has a horizontal tangent.\work{2.6cm}
  \item Write an equation of the line tangent to the graph of $f$ at $x=0$.\work{2.6cm}
\end{enumerate}
\end{frbox}

\newpage
\begin{solbox}
\begin{enumerate}[leftmargin=*,itemsep=4pt]
  \item $5e^{x}$.
  \item $2\cos x+3x^2$.
  \item $\dfrac4x-1$.
  \item $-\sin\dfrac\pi6=-\dfrac12$.
  \item The point is $\left(\dfrac\pi3,\sqrt3\right)$ and the slope is $2\cos\dfrac\pi3=1$: $y-\sqrt3=x-\dfrac\pi3$.
  \item $2-e^{x}=0$ at $x=\ln2$.
  \item The derivative of $e^{x}$ at 3: $e^{3}$.
  \item $\dfrac6x+2x=7$ gives $2x^2-7x+6=(2x-3)(x-2)=0$: $x=\dfrac32$ and $x=2$.
  \item \begin{enumerate}[label=(\alph*),leftmargin=*,itemsep=3pt]
    \item $f'(x)=4\cos x+2$.
    \item $\cos x=-\dfrac12$: $x=\dfrac{2\pi}{3}$ and $x=\dfrac{4\pi}{3}$.
    \item $f(0)=0$ and $f'(0)=6$: $y=6x$.
  \end{enumerate}
\end{enumerate}
\end{solbox}

\end{document}
